% TODO make hw stylesheet
\documentclass[10pt]{article}
\usepackage[letterpaper,top=1in]{geometry}
\usepackage[dvipsnames,svgnames]{xcolor}
\usepackage[shortlabels]{enumitem}
\usepackage[framemethod=TikZ]{mdframed}
\usepackage{amsmath,amssymb,amsthm}
\usepackage[colorlinks]{hyperref}
\usepackage{multicol}
\usepackage{tabularx}

\hypersetup{
	colorlinks=true,
	linkcolor=blue,
	filecolor=magenta,
	urlcolor=magenta,
	pdftitle={MAT 528 HW}
}

\usepackage{mathtools}
\usepackage{thmtools}
\usepackage[textsize=footnotesize]{todonotes}
\usepackage{titling}

\reversemarginpar
\mdfdefinestyle{mdbluebox}{roundcorner=10pt,innerbottommargin=9pt,
	linecolor=blue,backgroundcolor=TealBlue!5,}
\declaretheoremstyle[headfont=\sffamily\bfseries\color{MidnightBlue},
mdframed={style=mdbluebox},]{thmbluebox}

\mdfdefinestyle{mdbloobox}{frametitlefont=\bfseries,innerbottommargin=8pt,
	nobreak=true,backgroundcolor=TealBlue!5,linecolor=blue,}
\declaretheoremstyle[headfont=\bfseries\color{MidnightBlue},
mdframed={style=mdbloobox},headpunct={\\[3pt]},postheadspace=0pt,]{thmbloobox}

\mdfdefinestyle{mdredbox}{frametitlefont=\bfseries,innerbottommargin=8pt,
	nobreak=true,backgroundcolor=Salmon!5,linecolor=RawSienna,}
\declaretheoremstyle[headfont=\bfseries\color{RawSienna},
mdframed={style=mdredbox},headpunct={\\[3pt]},postheadspace=0pt,]{thmredbox}

\mdfdefinestyle{mdgreenbox}{linecolor=ForestGreen,backgroundcolor=ForestGreen!5,
	linewidth=2pt,rightline=false,leftline=true,topline=false,bottomline=false,}
\declaretheoremstyle[headfont=\bfseries\sffamily\color{ForestGreen!70!black},
mdframed={style=mdgreenbox},headpunct={ --- },]{thmgreenbox}

\mdfdefinestyle{mdorangebox}{linecolor=RedOrange,backgroundcolor=RedOrange!5,
	linewidth=2pt,rightline=false,leftline=true,topline=false,bottomline=false,}
\declaretheoremstyle[headfont=\bfseries\sffamily\color{RedOrange!70!black},
mdframed={style=mdorangebox},headpunct={ --- },]{thmorangebox}

\mdfdefinestyle{mdblackbox}{linecolor=black,backgroundcolor=RedViolet!5!gray!5,
	linewidth=3pt,rightline=false,leftline=true,topline=false,bottomline=false,}
\declaretheoremstyle[mdframed={style=mdblackbox}]{thmblackbox}

\declaretheoremstyle[spaceabove=6pt,spacebelow=6pt]{basehead}

\declaretheorem[style=basehead,name=Problem]{problem}
\declaretheorem[style=thmredbox,name=Problem,numbered=no]{problem*}
\declaretheorem[style=thmbloobox,name=Setup]{setup} % for config geo
\declaretheorem[style=thmredbox,name=Example,sibling=problem]{example}
\declaretheorem[style=thmredbox,name=$\bigstar$ Problem,sibling=problem]{niceprob}
\declaretheorem[style=thmbluebox,name=Theorem,sibling=problem]{theorem}
\declaretheorem[style=thmbluebox,name=Corollary,sibling=theorem]{corollary}
\declaretheorem[style=thmbluebox,name=Proposition,sibling=theorem]{prop}
\declaretheorem[style=thmbluebox,name=Lemma,numberwithin=problem]{lemma}
\declaretheorem[style=thmbluebox,name=Theorem,numbered=no]{theorem*}
\declaretheorem[style=thmbluebox,name=Lemma,numbered=no]{lemma*}
\declaretheorem[style=thmgreenbox,name=Claim,numberwithin=problem]{claim}
\declaretheorem[style=thmgreenbox,name=Claim,numbered=no]{claim*}
\declaretheorem[style=thmblackbox,name=Remark,numberwithin=problem]{remark}
\declaretheorem[style=thmblackbox,name=Remark,numbered=no]{remark*}
\declaretheorem[style=thmgreenbox,name=Definition,sibling=theorem]{definition}
\declaretheorem[style=thmgreenbox,name=Definition,numbered=no]{definition*}

\providecommand{\ol}{\overline}
\providecommand{\ceil}[1]{\left\lceil #1 \right\rceil}
\providecommand{\floor}[1]{\left\lfloor #1 \right\rfloor}
\newcommand{\dd}{\mathrm{d}}
\providecommand{\ul}{\underline}
\providecommand{\eps}{\varepsilon}
\providecommand{\half}{\frac{1}{2}}
\providecommand{\CC}{\mathbb C}
\providecommand{\FF}{\mathbb F}
\providecommand{\NN}{\mathbb N}
\providecommand{\QQ}{\mathbb Q}
\providecommand{\RR}{\mathbb R}
\providecommand{\ZZ}{\mathbb Z}
\providecommand{\dg}{^\circ}
\providecommand{\ii}{\item}
\providecommand{\setdiff}{\smallsetminus}
\providecommand{\alert}[1]{{\sffamily\textbf{\textcolor{blue}{#1}}}}
\providecommand{\inv}{^{-1}}
\providecommand{\mb}[1]{\mathbf{#1}}

\newenvironment{soln}{\begin{proof}[Solution]}{\end{proof}}

\providecommand{\pow}{\operatorname{Pow}}
\providecommand{\vocab}[1]{{\textbf{\textcolor{ForestGreen}{#1}}}}
\providecommand{\lcm}{\operatorname{lcm}}
\providecommand{\abs}[1]{\left|#1\right|}

\usepackage{mathrsfs}

% place title at top right
\setlength{\droptitle}{-8ex}
\pretitle{\begin{flushright}\Large\bfseries}
\posttitle{\par\end{flushright}}
\preauthor{\begin{flushright}}
\postauthor{\end{flushright}}
\predate{\begin{flushright}}
\postdate{\end{flushright}}

\title{MAT 528 HW04}
\author{\large Aditya Pahuja}
\date{\large Due 09/25}
\begin{document}
\maketitle

\begin{problem}
    [Munkres 22.1]
    Check the details of Example 3.
\end{problem}

\begin{soln}
    We are given the space $X = \{a,b,c\}$ with topology
    \begin{equation*}
	\{\varnothing, X, \{a\}, \{b\}, \{a,b\}\}
    \end{equation*}
    and the map $p\colon\RR\to X$ is given by
    \begin{equation*}
        p(x) = \begin{cases}
	    a & x>0\\
	    b & x<0\\
	    c & x=0.
        \end{cases}
    \end{equation*}
    $p$ is surjective because $p(1) = a$, $p(-1) = b$, and $p(0) = c$
    so each value in $X$ is achieved by $p$.
    To check continuity of $p$ we check that the pre-image of
    each open subset of $X$ is open in $\RR$:
    $p\inv(\varnothing) = \varnothing$ is open;
    $p\inv(X) = \RR$ is open;
    $p\inv(\{a\}) = (0,\infty)$ is open;
    $p\inv(\{b\}) = (-\infty,0)$ is open; and
    $p\inv(\{a,b\}) = (-\infty,0)\cup(0,\infty)$ is open.

    Finally, to check that $p$ is a quotient map we need to check that
    no non-open set $A\subseteq X$ has $p\inv(A)$ be an open subset of $\RR$.
    There are only three subsets of $X$ that aren't open,
    and we check them manually:
    $p\inv(\{c\}) = \{0\}$, $p\inv(\{a,c\}) = [0,\infty)$,
    and $p\inv(\{b,c\}) = (-\infty,0]$ are all not open in $\RR$
    because they do not contain any neighborhood of $0$.

    Thus, we have verified that $p\colon\RR\to X$ is a quotient map,
    which means that the topology above is the one induced by $p$.
\end{soln}

\begin{problem}
    [Munkres 22.2b]
    If $A\subseteq X$, a \vocab{retraction} of $X$ onto $A$ is
    a continuous map $r\colon X\to A$ such that $r(a) = a$ for each $a\in A$.
    Show that a retraction is a quotient map.
\end{problem}

\begin{soln}
    First, observe that $r$ is surjective because each $a\in A$
    is the image of the element $a\in X$ under $r$,
    since $r$ fixes points in $A$.

    Since $r$ is continuous, if $V\subseteq A$ is open then
    $r\inv(V)\subseteq X$ must be open by the definition of continuity.

    Conversely, suppose $r\inv(V)\subseteq X$ is open for some $V\subseteq A$.
    If $a\in r\inv(V)\cap A$, then $a = r(a)\in r(r\inv(V)) = V$
    where $a = r(a)$ because $a\in A$, so $r\inv(V)\cap A\subseteq V$.
    On the other hand if $a\in V$ then $a\in A$ because $V$ is a subset of $A$
    and $a\in r\inv(V)$ because $a = r(a)$, giving the reverse inclusion
    $V\subseteq r\inv(V)\cap A$. Therefore $V = r\inv(V)\cap A$,
    which means that $V$ is open in $A$.

    Thus $r$ is surjective and
    $V\subseteq A$ is open if and only if $r\inv(V)\subseteq X$ is open,
    which is what it means for $r$ to be a quotient map.
\end{soln}

\pagebreak

\begin{problem}
    [Munkres 22.4a]
    Define an equivalence relation on the plane $X = \RR^2$ as follows:
    \begin{equation*}
	(x_0,y_0)\sim(x_1,y_1)\text{ if }x_0+y_0^2 = x_1+y_1^2.
    \end{equation*}
    Let $X^\ast$ be the corresponding quotient space.
    It is homeomorphic to a familiar space; what is it?
\end{problem}

\begin{soln}
    Let $f\colon\RR^2\to\RR$ be the function $f(x,y) = x + y^2$.
    Then, $(x_0,y_0)\sim(x_1,y_1)$ means $f(x_0,y_0) = f(x_1,y_1)$,
    so elements of $X^\ast$ are exactly pre-images $f\inv(a)$ for $a\in\RR$;
    denote $f\inv(a)\in X^\ast$ by $[a]$.
    This means there is a bijection $h\colon X^\ast\to\RR$
    given by $h([a]) = a$, since there is exactly one equivalence class
    per $a\in\RR$ (and no equivalence class is empty
    because $f(a,0) = a+0^2 = a$ means $(a,0)\in[a]$).

    \begin{claim}
        $f$ is a quotient map.
    \end{claim}
    \begin{proof}
        $f$ is surjective because $f(r,0) = r$ for any $r\in\RR$.

	Let $\widetilde f\colon\RR^2\to\RR^2$ be the map
	$\widetilde f(x,y) = (f(x,y),0)$;
	it is a retraction of $X$ onto the $x$-axis
	because it is continuous
	(since the component functions $f$ and the zero function are) and
	$\widetilde f(r,0) = (f(r,0),0) = (r,0)$ for all $r\in\RR$.
	Thus by the previous problem it is a quotient map.

	Then, the map $\pi_1\colon\RR\times\{0\}\to\RR$ given by
	$\pi_1(x,0) = x$ is a homeomorphism, so
	$U\subseteq\RR$ is open if and only if
	$U\times\{0\}\subseteq\RR\times\{0\}$ is open,
	which in turn is open if and only if
	$\widetilde f\inv(U\times\{0\})\subseteq\RR^2$ is open,
	by definition of quotient map.
	Thus $f = \pi_1\circ\widetilde f\colon\RR^2\to\RR$ is a quotient map
	because $U\subseteq\RR$ is open if and only if
	$f\inv(U)\subseteq\RR^2$ is open.
    \end{proof}

    To finish, when $X^\ast$ has the quotient topology induced by $p$,
    we can use the fact that $p$ and $f$ are both quotient maps
    going out of $X$ to get that
    \begin{align*}
	U^\ast\subseteq X^\ast\text{ is open}
	&\iff p\inv(U^\ast) = \bigcup_{[a]\in U^\ast}[a]
	\subseteq X\text{ is open}\\
	&\iff f(p\inv(U^\ast)) = \{a\in\RR : [a]\in U^\ast\}
	= h(U^\ast) \subseteq\RR\text{ is open},
    \end{align*}
    This means that $h$ is a homeomorphism between $X^\ast$ and $\RR$,
    so the answer is $\boxed\RR$.
\end{soln}

\begin{problem}
    [Munkres 23.1]
    Let $\mathcal T$ and $\mathcal T'$ be two topologies on $X$.
    If $\mathcal T'\supseteq\mathcal T$, what does
    connectedness of $X$ in one topology imply about
    connectedness in the other?
\end{problem}

\begin{soln}
    If $X$ is connected with respect to $\mathcal T'$ then
    it must be connected with respect to $\mathcal T$.
    This is easier to see in the contrapositive form
    ``if $X$ is disconnected in $\mathcal T$ then
    it is disconnected in $\mathcal T'$'':
    if $A,B\in\mathcal T$ comprise a separation of $X$,
    then $A,B$ are also in $\mathcal T'$ so they comprise a separation of $X$
    with respect to the finer topology $\mathcal T'$.

    The converse is not true, though; for example, any set $X$
    with more than one element is connected in the indiscrete topology
    (because there are no nontrivial open sets)
    but disconnected in the discrete topology
    (because every subset is open, so if $A$ is a proper nonempty subset of $X$
    then $A\cup(X\setminus A)$ is a separation of $X$).
\end{soln}

\begin{problem}
    [Munkres 23.6]
    Let $A\subseteq X$. Show that if $C$ is a connected subspace of $X$
    that intersects both $A$ and $X\setminus A$,
    then $C$ intersects the boundary $\partial A$.
\end{problem}

\begin{soln}
    Recall that $\partial A = \ol A \setminus \operatorname{Int}A$,
    so $X$ is partitioned into the three sets
    $X\setminus\ol A$, $\partial A$, and $\operatorname{Int}A$
    (in particular, the three sets are pairwise disjoint).
    Suppose for the sake of contradiction that $C$ does not intersect $\partial A$.
    Let $C_1 = C\cap\operatorname{Int}A$ and $C_2 = C\cap(X\setminus\ol A)$.
    Since $\operatorname{Int}A$ and $X\setminus\ol A$ are open in $X$ and disjoint,
    $C_1$ and $C_2$ are open in $C$ and disjoint.
    Furthermore, $C_1$ and $C_2$ are both nonempty:
    if $C_1$ were empty, then $C = C_2$ would be contained in
    $X\setminus\ol A\subseteq X\setminus A$, so $C$ wouldn't intersect $A$,
    and similarly if $C_2$ were empty then $C$ would be contained in $A$
    so $C$ wouldn't intersect $X\setminus A$.
    Therefore, $C_1$ and $C_2$ are a separation of $C$,
    meaning $C$ is disconnected, which is a contradiction.
\end{soln}

\begin{problem}
    [Munkres 23.9]
    Let $A$ be a proper subset of $X$, and let $B$ be a proper subset of $Y$.
    If $X$ and $Y$ are connected, show that $(X\times Y)\setminus(A\times B)$
    is connected.
\end{problem}

\begin{soln}
    For any $b\in B$, the subspace $X\times\{b\}$ of $Z$
    is homeomorphic to $X$ via the homeomorphism $(x,b)\mapsto x$
    and for any $a\in A$, the subspace $\{a\}\times Y$ of $Z$
    is homeomorphic to $Y$ via the homeomorphism $(a,y)\mapsto y$,
    so in particular these subspaces are all connected.

    Fix $a_0\in A$. Then for any $b\in B$, the connected sets
    $X\times\{b\}$ and $\{a_0\}\times Y$ intersect at the point
    $\{(a_0,b)\}$, which implies that $T_b = (X\times\{b\})\cup(\{a_0\}\times Y)$
    is connected by Munkres's Theorem 23.3.
    Then, $\bigcap_{b\in B} T_b = \{a_0\}\times Y$ so the $T_b$
    all intersect at a common point, which means by Theorem 23.3 again
    that $\bigcup_{b\in B} T_b = (X\times B)\cup(\{a_0\}\times Y)$
    is connected.

    Similarly, fixing $b_0\in B$ and defining
    $T'_a = (\{a\}\times Y)\cup(X\times\{b_0\})$ for $a\in A$,
    we see that $\bigcap_{a\in A}T'_a$ is nonempty
    so $\bigcup_{a\in A}T'_a = (A\times Y)\cup(X\times\{b_0\})$ is connected.

    To finish, $(A\times Y)\cup(X\times\{b_0\})$ and
    $(X\times B)\cup(\{a_0\}\times Y)$ intersect at $(a_0,b_0)$,
    so their union is connected by Theorem 23.3.
    Since $X\times\{b_0\}\subseteq X\times B$ and $\{a_0\}\times Y\subseteq A\times Y$
    so their union
    \begin{equation*}
	\bigcup_{a\in A}T'_a\cup\bigcup_{b\in B}T_b =
	(X\times B)\cup(A\times Y) = (X\times Y)\setminus (A\times B)
    \end{equation*}
    is connected, which is what we wanted to show.
\end{soln}










\end{document}

